\section{Related Work}
\label{sec:related}

Neural operators learn mappings between function spaces and are therefore a natural model class for parametric PDE solution operators. DeepONet introduced a branch--trunk architecture for learning nonlinear operators from function samples and output coordinates \citep{lu2021deeponet}. Fourier Neural Operators (FNOs) parameterize global integral kernels in Fourier space and provide an efficient backbone for PDE families on regular grids \citep{li2021fno}. The general neural-operator framework formalizes these models as discretization-invariant compositions of integral operators and nonlinearities, with applications to Burgers' equation, Darcy flow, and Navier--Stokes \citep{kovachki2023neuraloperator}. Physics-Informed Neural Operators (PINOs) add PDE residual constraints to operator learning, combining supervised operator regression with physics-based optimization \citep{li2024pino}. Subsequent work extends neural operators to geometric settings: Geo-FNO learns deformations from irregular physical domains to latent uniform grids \citep{li2023geofno}, GINO combines graph and Fourier operator components for large-scale three-dimensional PDEs on varying geometries \citep{li2023gino}, and NORM uses Laplacian eigenfunctions to define neural operators on Riemannian manifolds \citep{chen2024norm}. Operator-learning benchmarks such as PDEBench and The Well provide reusable datasets and baselines for time-dependent PDEs and broader physical simulation families \citep{takamoto2022pdebench,ohana2024well}. This literature establishes FNO-style models as computationally efficient solution-operator backbones, while also showing that domain geometry, discontinuities, and discretization effects require additional structure \citep{lanthaler2023discontinuous}.

Symmetry has been used as an inductive bias through architectures, losses, and data transformations. Group-equivariant convolutional networks replace ordinary convolution by group convolution to share weights across symmetry orbits \citep{cohen2016gcnn}; steerable CNNs generalize this principle through representation-constrained kernels \citep{weiler2019e2cnn}; and gauge-equivariant CNNs extend equivariance from global transformations to local frame changes on manifolds \citep{cohen2019gauge}. For operator learning, Group-Equivariant FNOs encode translation, rotation, and reflection symmetries in Fourier layers \citep{helwig2023gfno}. A parallel line derives symmetries from the governing differential equations. Lie Point Symmetry Data Augmentation constructs solution-preserving transformations from a PDE's Lie point symmetry group and uses them to reduce sample complexity for neural PDE solvers \citep{brandstetter2022lpsda}. Lie-symmetry losses for PINNs enforce conservation of solution families under infinitesimal generators \citep{akhoundsadegh2023liepinn}. Generalized Lie symmetries for PINOs further argue that evolutionary representatives can provide stronger training signals than point symmetries alone \citep{wang2025generalizedliepino}. These works support the central premise that PDE symmetries should constrain the learned solution operator, not only the training data distribution.

Model-agnostic symmetry methods provide the closest precedent for the proposed direction. Tangent propagation used transformation tangents to impose local invariance constraints on neural networks \citep{simard1991tangentprop}. Lie Algebra Canonicalization uses infinitesimal generators to construct canonicalized inputs for pretrained models and neural operators under arbitrary Lie groups, including non-compact groups commonly arising in PDEs \citep{shumaylov2025lielac}. PACE-FNO applies a related canonicalization strategy to FNOs under symmetry-induced shifts, estimating a Lie-algebraic input frame, predicting in a canonical frame, and restoring the output frame for translations, rotations, and Galilean boosts \citep{xu2026pacefno}. LOCO occupies a distinct point in this design space. Unlike supervised Lie augmentation, its consistency term needs no transformed label; unlike hard-equivariant layers, it does not restrict the backbone; unlike canonicalization, it does not alter inference; and unlike tangent propagation, it uses ordinary evaluations at finite group elements rather than an explicit Jacobian--vector product at the identity. Relative to a plain finite-transform consistency loss, the local \(\epsilon^2+\eta\) normalization removes the leading quadratic dependence of the loss scale on step magnitude. The revision tests the finite-versus-tangent distinction and reports a scale-matched unnormalized control rather than treating these objectives as interchangeable.
